\section{The case \texorpdfstring{$\ell=0$}{l=0} is a CSP on the intersection graph}\label{sec:l0}
\begin{observation}\label{obs:one}
Let $S$ be a solution with $|S|\le|H|$. Then $S$ has exactly one pair in each $\pairs(P_i)$ and no other pair, and $G\sd S$ is a cluster graph iff no vertex triple is an induced $P_3$ of $G\sd S$. Consequently $(G,H,0)$ is a YES instance iff the following CSP is satisfiable: variables $x_i\in\{1,2,3\}$ (which pair of $P_i$ is flipped), and for every vertex triple $t$ the constraint ``$t$ is not an induced $P_3$ after flipping'', whose truth depends on the owners' variables of the packed pairs of $t$ only (unpacked pairs keep their status in $G$).
\end{observation}
\begin{proof}
By Lemma~\ref{lem:decomp}, $|x_i|\ge1$ for all $i$ and $|S|=\sum|x_i|+|U|\le|H|$ force $|x_i|=1$ and $U=\emptyset$. The rest is the definition of a cluster graph.
\end{proof}
The CSP has $\Theta(n^3)$ constraints, not $O(|H|)$.
\begin{observation}\label{obs:two}
The set of owners of the packed pairs of any triple is a clique of $I(H)$ (at most three nodes). Hence the primal graph of the CSP is a subgraph of $I(H)$ (not every edge of $I(H)$ is witnessed by a triple), and a tree decomposition of $I(H)$ is one of the primal graph.
\end{observation}
\begin{proof}
Two pairs of one triple share a vertex $v$; their owners both contain $v$, so they coincide or are adjacent in $I(H)$.
\end{proof}
\begin{theorem}\label{thm:l0}
Given $(G,H)$ and a tree decomposition of $I(H)$ of width $w$, CEaMP with $\ell=0$ is solvable in $3^{w+1}n^{O(1)}$ time.
\end{theorem}
Theorem~\ref{thm:l0} is the case $\ell=0$ of Theorem~\ref{thm:A} below (with a larger base there, $7$, because for $\ell>0$ a packed $P_3$ may receive several edits); its proof is the same dynamic program with domain $\{1,2,3\}$. Observation~\ref{obs:two} is what makes the interaction of disjoint $P_3$'s at $\ell=0$ completely visible in $I(H)$: no transitivity-mediated interaction exists beyond shared vertices.
